Even when these reduce to the Chern classes of locally free Modules (@@M0@@ a scheme with an ample invertible Module, or a compact space), these Chern classes ``with supports'' have not been constructed, and their construction will doubtless have to proceed together with that of the general Chern classes postulated in 7.1\SGAAuthorFootnote{*}{Added in October 1968. The problems raised here have been solved in the setting of \uline{Hodge cohomology} by L. ILLUSIE; cf. [13]}

It is probable that Chern classes with supports will have a role to play in proving the cohomological Riemann--Roch formula without denominators (cf. 3.5). On the other hand, they should allow an improved version of the Riemann--Roch formula (with or without denominators) by introducing families of supports into that formula.

\sgaStableTarget{sga6.e14.sec.8}\section*{8. Chow ring of regular schemes}

The problem is to define rational equivalence of cycles and to study the resulting ring of classes (in particular its functorial properties) for a regular noetherian scheme @@M1@@ admitting an ample invertible Module, generalizing the theory known when @@M2@@ is smooth over a field [4]. The key technical result will again be a ``moving lemma'', ensuring that a cycle can be varied within its class so as to intersect a given cycle in the correct dimension.

The problem considered is quite natural a priori, particularly since we have seen (4.5 to 4.8) that the substitutes @@M3@@, @@M4@@ for the Chow ring @@M5@@ are hardly satisfactory beyond a theory ``modulo torsion''. When @@M6@@ is obtained from a smooth projective algebraic scheme @@M7@@ over a field @@M8@@ by taking the local scheme at the vertex of a projecting cone @@M9@@ and removing the origin, one finds (admitting that @@M10@@ is indeed the direct limit of the Chow rings of the @@M11@@, where @@M12@@ ranges over the open neighbourhoods in @@M13@@ of the vertex @@M14@@) that @@M15@@ is canonically isomorphic to the quotient ring @@M16@@, where @@M17@@ is induced by the projective immersion of @@M18@@ under consideration. This shows that knowledge of the Chow ring @@M19@@ (of a nonclassical type) is equivalent to knowledge of @@M20@@, which might be called the ``primitive part'' of the Chow ring @@M21@@ (in a sense suggested by the Lefschetz--Hodge cohomological theory of nonsingular projective algebraic varieties). This example shows that the ``local Chow rings'' of local rings with a singularity
